\section{Conclusion}
\label{sec:conclusion}

We began with a counterintuitive claim: that optimizing a language model for alignment metrics makes it \emph{less} aligned with actual human judgment---and that we can quantify exactly how fast. We have demonstrated that this is not merely a theoretical possibility but a measurable, replicable empirical property of RLHF optimization.

The calibration-alignment divergence gap grows at $\beta \approx 0.143$--$0.160\ \text{nat}^{-1}$ of KL optimization budget---high enough to traverse more than half the normalized scale within 4--5 nats of KL beyond the divergence onset point, and consistent enough to replicate within 12\% across two independent model families.

Our four contributions: (1) the calibration-alignment divergence construct (named and quantified); (2) slopes with confidence---$\beta = 0.143\ \text{nat}^{-1}$ ($R^2 = 0.958$, $p < 10^{-6}$) with both CIs strictly positive; (3) cross-dataset replication with slope ratio 1.116; (4) the normalized divergence gap metric as a cross-study comparison instrument.

Future directions grounded in our results: (a) raw data access from original authors for exact $\beta$ with digitization uncertainty propagation; (b) piecewise linear regression on Gao data to characterize the pre-onset regime; (c) behavioral study linking evaluation calibration divergence to user over-reliance; (d) coverage ratio $R$ computation across the 400-paper bidirectional alignment survey corpus; (e) cross-scale replication at multiple RM sizes to test whether $\beta \approx 0.14$--$0.16\ \text{nat}^{-1}$ is scale-invariant.

We demonstrated that $\beta \approx 0.143\ \text{nat}^{-1}$---quantifying exactly how fast RLHF optimization degrades evaluation calibration to human judgment. A field that measures alignment by proxy scores alone, without tracking held-out human preference, systematically overestimates how aligned its models are as optimization proceeds.
